\paragraph{Conventions.} Belief time is the report day; two reports on the same day are ordered by id, and a query at belief time $\tau$ sees every report with $t_{\mathrm{rep}}\le\tau$. A report with no temporal cue is anchored at its report day; \texttt{since}~$t$ anchors it at $t$ (which may precede the report day); \texttt{at}~$t$ asserts the value at $t$ only. A query before the first anchor on a key is \texttt{unknown}; a query inside a gap between two different accepted anchors on a single-valued key has both neighbours as candidates (unknown change point); a reversal is two changes. Withdrawn evidence is never restored. A correction whose target is absent from the stream is an ordinary assertion; a correction of a withdrawn report is accepted and has no further effect; a same-origin correction later withdrawn leaves its target withdrawn (withdrawal is not undone); a source withdrawal removes every report of that source, including reports it makes after the withdrawal, for all later belief times. Multi-valued \texttt{competing\_values=false} attributes never dispute; their runs end only through withdrawal or correction.

\paragraph{Decision table: unknown, empty, false.} The three conventions of \S\ref{sec:semantics} applied to the cases a reader asked about (verified on toy streams in \texttt{tests/}):

\begin{center}\small\begin{tabular}{lccc}\toprule
Situation at the query & single-valued slot & multi-valued slot & yes/no ``holds $v$'' \\\midrule
no report on the key & \texttt{unknown} & \texttt{established} $\emptyset$ & \texttt{established} false \\
every report on the key withdrawn & \texttt{unknown} & \texttt{established} $\emptyset$ & \texttt{established} false \\
belief time before the first report & \texttt{unknown} & \texttt{established} $\emptyset$ & \texttt{established} false \\
undeclared entity or attribute & \texttt{unknown} & \texttt{unknown} & \texttt{unknown} \\
ladder exhausted (every report rejected) & \texttt{unknown} & \texttt{unknown} & \texttt{unknown} \\
$v$ never asserted, key otherwise established & \texttt{established} $\neq v$ & set without $v$ & \texttt{established} false \\\bottomrule
\end{tabular}\end{center}
The first three rows have exactly one admissible interpretation with an empty timeline, which yields no candidate for a single-valued slot (hence \texttt{unknown}), the empty set for a multi-valued slot, and falsity for every ``holds'' proposition (closed world over admitted evidence); the undeclared and ladder-exhausted rows have no interpretation, hence \texttt{unknown} throughout. The difference between a missing key and an exhausted ladder is reported through the ladder counter, not through the status.

\paragraph{Relaxation ladder (totality).} The per-key enumeration can leave a key with no admissible interpretation on extracted or adversarial streams (e.g.\ a change presupposition whose only witness is withdrawn, or under P1). The oracle is made total by a declared ladder applied per key only when the admissible set is empty: (1) drop A-CHG presuppositions; (2) additionally fall back to literal P0 dispute (no cue shielding); (3) reject every report on the key (status \texttt{unknown}). Each application is counted; the ladder never fired on the 500 generated Setting 1 streams under P0c, on the 100 extracted Setting 2 streams, or on the 30 Setting 3 streams (counted on 3 October 2026).

\paragraph{Rules and the per-key factorisation.} Derived attributes are forward-chained per interpretation; the released domains are acyclic with derivation depth $\le2$, and cyclic rule sets are outside the declared scope. A defeasible rule contributes nothing under an interpretation in which its exception holds (explicit truth of the exception, not absence). A derived value with several independent derivations keeps its candidate when one premise is withdrawn as long as another derivation survives; provenance records every derivation path that yields a candidate. The oracle computes interpretation sets \emph{per key} and lets the rule engine branch over the per-key candidate unions. This is exact, not an approximation, under two conditions that the declared semantics and rule sets satisfy: (i) every admissibility condition (A-ERR, A-CORR, A-CHG, A-CUE, A-CONS, A-SU) relates reports on the same key, and a correction that refers to a report on another key acts only through admission (A-SELF withdrawal), so the set of global interpretations is the Cartesian product of the per-key sets; (ii) no key occurs more than once on a derivation path, so a derived candidate obtained from one candidate per key is realised by the global interpretation that picks those per-key interpretations. An earlier draft of this appendix said that the factorisation ``can over-generate alternatives'' in general; the correct statement is that it over-generates exactly when condition (ii) fails. The reviewer's diagnostic shows it: with $U\leftarrow X$, $V\leftarrow X$ and $W\leftarrow U\wedge V$ and $X$ unresolved between $A$ and $B$, the global semantics gives $W\in\{A,B\}$ while the per-key procedure also produces the impossible mixed derivation (reported as a ``no value'' alternative); the released check reproduces this (\texttt{oracle\_validation.py --counterexample}) and a static test on the rule sets (the attribute closures of a rule's body atoms must be pairwise disjoint) flags that domain and passes every released domain. Independent alternative proofs of one derived value do not occur in the released rule sets (one rule per derived attribute), and the validation below covers defeasible exceptions and multi-valued alternatives through the generator's own domain. Validation (\texttt{scripts/oracle\_validation.py}): on \ovStreams{} streams (the three hand-written references and 300 small generated ones, seeds 40000--) and every value query whose dependency closure admits at most \ovMaxEnum{} reports (\ovSkipped{} larger closures skipped), the per-key gold was compared with a brute-force \emph{global} computation that enumerates every TRUE/ERR and CONT/ENDED labelling of the closure's reports and derives each query under each global interpretation's own timelines for all keys at once: \ovQueries{} value queries (\ovDerived{} on derived slots), \ovDisagree{} disagreements in status, value or alternative set; the product structure $|\text{global}|=\prod_{\text{keys}}|\text{per-key}|$ held in all \ovQueries{} checks, and the per-key gold on the closure sub-stream equalled that on the full stream in every case (\ovClosureFail{} failures), which is the independence of keys outside the closure.

\paragraph{Worked example, revisited (reviewer question).} With \emph{residence} multi-valued and changeable, after ``Alex lives in London'' (o1) and ``Alex moved from London to Paris'' (o2), A-CHG requires some earlier London report to be accepted, A-CUE shields o2 from o1, and the \textsc{ended} label on London's run is \emph{forced} by the \emph{from} clause of the change cue (\S\ref{sec:semantics}, ``a change-from cue forces \textsc{ended} on the run of the previous value''); the written clause of the earlier draft omitted this sentence and the implementation has always had it. Without the \emph{from} clause (``Alex has moved to Paris'') London's continuation stays open, and ``Did Alex move?'' is \texttt{established} while ``Does Alex still live in London?'' is \texttt{possible}.
